\chapter{Boolean Analysis — Encoding}
%%% generated by the blueprint pipeline from TCSlib.BooleanAnalysis.Switching.Encoding
\section{Overview}

This module defines the Razborov encoding and decoding maps used in the proof of
the switching lemma. The encoder turns a ``bad'' restriction $\rho$ into a pair
$(\gamma, \mathrm{aux})$, where $\gamma$ extends $\rho$ by fixing the variables
along the deepest path of the canonical decision tree for $f|_\rho$ to their
satisfying directions, and $\mathrm{aux}$ records, clause by clause, the data needed
to invert the process; the decoder recovers $\rho$ from $(\gamma, \mathrm{aux})$.

%%%%%%%%%%%%%%%%%%%%%%%%%%%%%%%%%%%%%%%%%%%%%%%%%%%%%%%%%%%%%%%%%%%%%%%%%%%%%%%
\section{Declarations}

\begin{definition}[Free literals of a term under a restriction]
\lean{SwitchingLemma2.Term.freeLiterals}
\leanok
\uses{Literal, SwitchingLemma2.Restriction, SwitchingLemma2.Restriction.freeVars, Term}
Given a term $t$ (a conjunction of literals) and a restriction $\rho$, this is the list
of literals of $t$ whose underlying variable is still free under $\rho$, i.e. the
sublist of $t$ obtained by keeping those literals $l$ with $l.\mathrm{var} \in \rho.\mathrm{freeVars}$.
\end{definition}

\begin{definition}[Process one clause's free literals against the path]
\lean{SwitchingLemma2.processClauseLits}
\leanok
\uses{Literal, SwitchingLemma2.Restriction}
Consumes the free literals of a single clause (each paired with its position index)
against the canonical decision-tree path, one path entry per literal. For each free
literal it fixes the literal's variable to its satisfying direction in $\sigma$,
records the path direction in $\rho_0$ (mirroring the canonical decision-tree
branching), and appends the pair (literal position in clause, path direction) to the
auxiliary data. It returns the remaining path together with the updated $\rho_0$,
$\sigma$, and the accumulated clause auxiliary data.
\end{definition}

\begin{definition}[Razborov encoding]
\lean{SwitchingLemma2.razborovEncode}
\leanok
\uses{DNF, DecisionTree.deepPath, SwitchingLemma2.Restriction, SwitchingLemma2.canonicalDTree}
The Razborov encoding of a DNF $f$ and a restriction $\rho$ with parameters $w, d$.
It takes the first $d$ steps of the deepest path of the canonical decision tree for
$f|_\rho$, then repeatedly selects the first clause of $f$ not killed by the current
restriction, processes all of that clause's free literals against the path, and emits
a termination marker $(w, \mathrm{false})$. The result is a pair $(\gamma, \mathrm{aux})$,
where $\gamma$ extends $\rho$ by fixing $d$ variables to their satisfying directions
and $\mathrm{aux}$ consists of per-clause blocks separated by the termination markers.
\end{definition}

\begin{definition}[Razborov decoding]
\lean{SwitchingLemma2.razborovDecode}
\leanok
\uses{DNF, SwitchingLemma2.Restriction}
The inverse of the Razborov encoding: given the DNF $f$, parameter $w$, and a pair
$(\gamma, \mathrm{aux})$ produced by the encoder, it recovers the original restriction
$\rho$. It repeatedly finds the first clause not killed by the current restriction,
processes that clause's auxiliary block (unfixing each recorded variable in $\sigma$
and recording its path direction in $\rho_0$ until a termination marker is reached),
and returns the resulting restriction.
\end{definition}

\begin{lemma}[Processing does not lengthen the path]
\lean{SwitchingLemma2.processClauseLits_path_le}
\leanok
\uses{Literal, SwitchingLemma2.Restriction, SwitchingLemma2.processClauseLits}
The remaining path returned by \lean{processClauseLits} is no longer than the input
path: the length of the first component of $\mathrm{processClauseLits}\,(\mathit{lits},
\mathit{path}, \rho_0, \sigma)$ is at most the length of $\mathit{path}$.
\end{lemma}
